\section{Metodologia}

\subsection{Formulação do Problema}
\begin{frame}{Formulação do Problema}
  Descreva aqui a formulação formal do problema abordado.
  
  \begin{block}{Definição Formal}
    Apresente a definição matemática ou conceitual do problema:
    \begin{itemize}
      \item $x$: Descrição da variável de decisão.
      \item $f(x)$: Função objetivo a ser otimizada.
      \item $g(x) \leq 0$: Restrições do problema.
    \end{itemize}
  \end{block}
  
  \begin{itemize}
    \item Explique brevemente as premissas e simplificações adotadas.
    \item Justifique as escolhas de modelagem.
  \end{itemize}
\end{frame}

\subsection{Abordagem Proposta}
\begin{frame}{Abordagem Proposta}
  Apresente a abordagem ou algoritmo utilizado para resolver o problema.

  \begin{exampleblock}{Método Proposto}
    Descreva em alto nível o método ou algoritmo empregado, suas etapas principais e por que foi escolhido.
  \end{exampleblock}
  
  \begin{itemize}
    \setlength{\itemsep}{10pt}
    \item \textbf{Etapa 1}: Descrição da primeira etapa do método.
    \item \textbf{Etapa 2}: Descrição da segunda etapa do método.
    \item \textbf{Etapa 3}: Descrição da terceira etapa do método.
  \end{itemize}
\end{frame}

\subsection{Detalhes Técnicos}
\begin{frame}{Detalhes Técnicos}
  \small
  \begin{itemize}
    \setlength{\itemsep}{10pt}
    \item \textbf{Estruturas de Dados}: Descreva as estruturas de dados utilizadas na implementação.
    \item \textbf{Complexidade Computacional}: Apresente a análise de complexidade do algoritmo.
    \item \textbf{Ferramentas e Linguagens}: Liste as tecnologias utilizadas na implementação.
  \end{itemize}
  
  \begin{table}[h]
    \centering
    \small
    \begin{tabular}{llr}
      \toprule
      \textbf{Componente} & \textbf{Descrição} & \textbf{Complexidade} \\
      \midrule
      Etapa 1 & Descrição breve & $O(n)$ \\
      Etapa 2 & Descrição breve & $O(n \log n)$ \\
      Etapa 3 & Descrição breve & $O(n^2)$ \\
      \bottomrule
    \end{tabular}
    \caption{Resumo de complexidade por etapa}
  \end{table}
\end{frame}


\section{Resultados}

\subsection{Experimentos}
\begin{frame}{Configuração Experimental}
  \begin{columns}[t]
    \begin{column}{0.5\textwidth}
      \begin{block}{Ambiente}
        \begin{itemize}
          \scriptsize
          \item \textbf{Hardware}: Descrição do hardware utilizado.
          \item \textbf{Software}: Linguagem, bibliotecas e versões.
          \item \textbf{Datasets}: Bases de dados ou instâncias utilizadas.
        \end{itemize}
      \end{block}
    \end{column}
    \begin{column}{0.5\textwidth}
      \begin{block}{Métricas de Avaliação}
        \begin{itemize}
          \scriptsize
          \item \textbf{Métrica 1}: Descrição da métrica.
          \item \textbf{Métrica 2}: Descrição da métrica.
          \item \textbf{Métrica 3}: Descrição da métrica.
        \end{itemize}
      \end{block}
    \end{column}
  \end{columns}
\end{frame}

\subsection{Resultados Obtidos}
\begin{frame}{Resultados Obtidos}
  \begin{columns}[t]
    \begin{column}{0.54\textwidth}
      \begin{table}
        \centering
        \small
        \begin{tabular}{lrr}
          \toprule
          \textbf{Método} & \textbf{Métrica 1} & \textbf{Métrica 2} \\
          \midrule
          Baseline      & $0.75$ & $1.23$ \\
          Proposto      & $\mathbf{0.92}$ & $\mathbf{0.87}$ \\
          Alternativo   & $0.81$ & $1.05$ \\
          \bottomrule
        \end{tabular}
        \caption{Comparação de métodos}
      \end{table}
    \end{column}
    \begin{column}{0.46\textwidth}
      \begin{block}{Síntese dos Resultados}
        \begin{itemize}
          \scriptsize
          \item \textbf{Resultado principal}: Descreva o principal achado.
          \item \textbf{Ganho observado}: Quantifique a melhoria.
          \item \textbf{Observações}: Comentários adicionais.
        \end{itemize}
      \end{block}
    \end{column}
  \end{columns}
\end{frame}

\subsection{Discussão}
\begin{frame}{Discussão e Análise}
  \begin{itemize}
    \setlength{\itemsep}{10pt}
    \item \textbf{Interpretação dos resultados}: Analise os resultados obtidos e o que eles significam no contexto do problema.
    \item \textbf{Comparação com a literatura}: Como seus resultados se comparam com os trabalhos relacionados?
    \item \textbf{Pontos fortes}: Destaque os pontos fortes da abordagem proposta.
    \item \textbf{Limitações observadas}: Quais limitações foram identificadas durante os experimentos?
  \end{itemize}
\end{frame}
